\documentclass[Total.tex]{subfiles}

\begin{document}
	\chapter{Vector Bundles}
		\section{Quasi-Vector Bundles}
			
			Let $k=\RR,\CC$
			\begin{Def}
				A \textbf{quasi-vector bundle} with base $X$ is given by:
				\begin{itemize}
					\item A finite dimensional $k$-vector space $E_x$ for every point $x\in X$;
					\item A topology $E=\coprod E_x$;
				\end{itemize}
				
				A quasi-vector bundle 
			\end{Def}
			
			\textbf{Remark} This also is called family of vectors.
			
			\begin{example}
				内容...
			\end{example}
			
			\begin{example}
				内容...
			\end{example}
	
		\section{Vector Bundles}
			
			\begin{Def}
				An $n$-dimensional vector bundle is a map $p:E\rightarrow B$ together with 
				
				\begin{itemize}
					\item a real vector space structure on $p^{-1}(b)$(\textbf{fibre}) for each $b\in B$;
					
					\item There exists a covering of $B$ by $U_\alpha$ such that: Equipped with homeomorphism
					
					\begin{gather*}
						h_\alpha:p^{-1}(U_\alpha)\stackrel{\cong}{\rightarrow} U_\alpha\times \RR^n
					\end{gather*}
					
					which is a homemorphism. And $h_{\alpha}|p^{-1}(b):p^{-1}(b)\rightarrow \{b\}\times K^n$ is a vector isomrophism.
					
					such $h_\alpha$ is called \textbf{trivialization} of the vector bundle. 
				\end{itemize}
			\end{Def}
			
			In principle, $K=\CC,\RR$. 
			
			\begin{example}
				内容...
			\end{example}
			
			\begin{example}
				内容...
			\end{example}
			
			\begin{example}
				(Tangent Bundle)
			\end{example}
			
			\begin{example}
				(Normal Bundle)
			\end{example}
			
			\begin{example}
				内容...
			\end{example}
			
			\subsection{title}
		
		\section{Classifying Vector Bundles}
			\subsection{Pullback Bundles}
				
			\subsection{Clutching Functions}
			
			\subsection{Cell Structures on Grassmannians}
		\section{Homotopy Theory of Vector Bundles}
		
			\begin{theorem}
				Let $X$ be a compact space and $E$ be a vector bundle over $X\times I,I=[0,1]$. Let
				\begin{gather*}
					\alpha_t:X\rightarrow X\times I\qquad \prod:X\times I\rightarrow X\\
					x\mapsto (x,t)\qquad (x,u)\mapsto x
				\end{gather*}
				Then the vector bundles $E_0=\alpha_0^*(E)$ is isomorphic to $E_1=\alpha_1^*(E)$.
			\end{theorem}
			\begin{proof}
				内容...
			\end{proof}
			
			\begin{theorem}
				内容...
			\end{theorem}
			\begin{proof}
				内容...
			\end{proof}
			
			\begin{theorem}
				内容...
			\end{theorem}
			\begin{proof}
				内容...
			\end{proof}
			
		\section{Metrics and Forms on Vector Bundles}
	\chapter{K-Theory}
		\section{The Functor $K(X)$}
	\chapter{Operations}
	\chapter{Bott Periodity}
	\chapter{Computation of Some K-Groups}
	\chapter{Applications of K-Theory}
	
	\chapter{Stiefel-Whitney Classses}
	\chapter{Cell-Structure for Grassmannian Manifolds}
	\chapter{The Cohomology Ring}
	\chapter{Chern-Weil Theory}
			
\end{document}